%Derive angular momentum and total energy as constants of the motion
\section{Introduction}
A central force is a force that acts only in the radial direction. In other words the force only depends on the distance from the centre. As such a central force can be defined like so
\begin{equation}
\mathbf{F} = f(r) \hat{\mathbf{r}} \label{eq:central1}
\end{equation}

The equation of motion for a central force is then
\begin{equation}
\ddot{\mathbf{r}} = \dot{\mathbf{v}}= \frac{\mathbf{F}}{m}= \frac{f(r)}{m} \hat{\mathbf{r}} \label{eq:central2}
\end{equation}
where we have used $\mathbf{v}=\dot{\mathbf{r}}$.

\section{Angular Momentum}

One important parameter for a central force is the angular momentum. The angular momentum is defined as the cross product of the position vector $\mathbf{r}$ and the momentum vector $\mathbf{p}=m\mathbf{v}$

\begin{equation}
\mathbf{L}=\mathbf{r}\times\mathbf{p} = m \mathbf{r}\times\mathbf{v}
\end{equation}

The mass $m$ is fixed so it is more convenient to use the specific angular momentum
\begin{equation}
\mathbf{h}=\frac{\mathbf{L}}{m}= \mathbf{r}\times\mathbf{v}
\end{equation}

\subsection{Conservation of Angular Momentum}

We can show that angular momentum is a constant of the motion by showing it's rate of change with respect to time is zero, ie it is fixed and unchanging. We start then by taking the derivative with respect to time.

\begin{align}
\dot{\mathbf{h}}&= \frac{d}{dt}(\mathbf{r}\times\mathbf{v}) \nonumber \\
&= \dot{\mathbf{r}}\times\mathbf{v}+\mathbf{r}\times\dot{\mathbf{v}} \nonumber \\
&= \mathbf{v}\times\mathbf{v}+\mathbf{r}\times\dot{\mathbf{v}}\nonumber \\
&= \mathbf{r}\times\dot{\mathbf{v}}
\end{align}
We have used the fact that the vector product of any vector with itself is zero. We can now substitute in equation \refpage{eq:central2} for $\dot{\mathbf{v}}$
\begin{align}
\dot{\mathbf{h}}&= \mathbf{r}\times\dot{\mathbf{v}}\nonumber \\
&= \mathbf{r} \times  \frac{f(r)}{m} \hat{\mathbf{r}}\nonumber \\
&= \frac{f(r)}{mr} \mathbf{r} \times \mathbf{r}\nonumber \\
&= \mathbf{0}
\end{align}
This shows that the angular momentum is a constant of the motion for a central force.

\section{Potential Energy}

The potential energy of a particle at a point is the amount of work that an external agent must do against the force exerted by the system on that particle from a chose reference point where we define the potential energy to be zero to that point.

In the case of a central force it is typical to choose the reference point to be at a radius of $\infty$ for attractive forces and to be at a radius of zero for repulsive forces. Additionally note that it is the work done against the force applied to the particle by the system, so the force the external agent must apply is equal and opposite the force applied by the system. This means when calculating the work done by the external agent, the force must be the negative of the force applied by the system. 

Here we want the potential energy at a radus r, so we integrate from $\infty$ to $r$
\begin{align}
V(r)&=\int_{\infty}^{r} -\mathbf{F}(r') \cdot d\mathbf{r}'  \nonumber \\
&=-\int_{\infty}^{r} f(r') \hat{\mathbf{r}} \cdot dr' \hat{\mathbf{r}} \nonumber \\
&=-\int_{\infty}^{r} f(r') dr'\label{eq:central-potentialEnergy}
\end{align}
If we take the derivative of both sides with respect to $r$ we get
\begin{align}
\frac{d V}{dr}&=-f(r)\label{eq:central3}
\end{align}

\section{Kinetic Energy}
The kinetic energy of a particle is given by
\begin{equation}
T(v)=\frac{1}{2}m\mathbf{v}\cdot\mathbf{v} = \frac{1}{2} m v^2\label{eq:central-KineticEnergy}
\end{equation}

\section{Total Energy}
The total energy is the sum of kinetic energy and potential energy
\begin{equation}
E =\frac{1}{2} m v^2 + V(r)
\end{equation}

Again it is often easier to use the specific total energy $\epsilon$

\begin{equation}
\epsilon=\frac{E}{m}=\frac{1}{2} v^2 + \frac{V(r)}{m}
\end{equation}
to make this a little easier still we will let $\tilde{V}=\frac{V(r)}{m}$, so we get
\begin{equation}
\epsilon=\frac{E}{m}=\frac{1}{2} v^2 +\tilde{V}
\end{equation}

\subsection{Conservation of Total Energy}
As before with the angular momentum, we can show that the total energy is a constant of motion by taking the derivative with respect to time

\begin{align}
\frac{d\epsilon}{dt}&=\frac{d}{dt}\left(\frac{1}{2} v^2 + \tilde{V}\right)\nonumber \\
&=\frac{1}{2} \frac{d v^2}{dt}+ \frac{d\tilde{V}}{dt}\nonumber \\
&=\frac{1}{2} \frac{d v^2}{dv}\frac{d v}{dt}+ \frac{d\tilde{V}}{dr}\frac{d r}{dt}\nonumber \\
&=\frac{1}{2} 2v\frac{d v}{dt}-\frac{f(r)}{m}v\nonumber \\
&=v\left(\dot{v}-\frac{f(r)}{m}\right)\label{eq:central4}
\end{align}
From equation \refpage{eq:central2}
\begin{align}
\dot{v} &= \sqrt{\dot{\mathbf{v}}\cdot\dot{\mathbf{v}}}\nonumber \\
&=\sqrt{ \frac{f(r)}{m} \hat{\mathbf{r}} \cdot \frac{f(r)}{m} \hat{\mathbf{r}}}\nonumber \\
&= \frac{f(r)}{m} \sqrt{ \hat{\mathbf{r}} \cdot \hat{\mathbf{r}}}\nonumber \\
&= \frac{f(r)}{m} \label{eq:central5}
\end{align}
we can then substitute equation \refpage{eq:central5} into equation \refpage{eq:central4} giving

\begin{align}
\frac{d\epsilon}{dt}&=v\left(\dot{v}-\frac{f(r)}{m}\right)\nonumber \\
&=v\left(\frac{f(r)}{m}-\frac{f(r)}{m}\right)\nonumber \\
&= 0
\end{align}
This then shows that total energy is a constant of the motion for a central force

\section{Summary}

In this chapter we have shown that both the angluar momentum and the total energy is conserved under a central force.